Un chimiste a préparé au laboratoire une quantité de cristaux d'acide benzoïque
de masse $m_{0}=\qty{244}{\mg}$. Après l'avoir dissout totalement dans de l'eau
distillée, il a obtenu une solution aqueuse ($S_{0}$) de
volume $V_{0}=\qty{100}{\mL}$ et de $\mathrm{pH}\simeq 2,95$.
\begin{enumerate}
    \item Écrire l'équation de la réaction modélisant la transformation ayant
          lieu entre l'acide benzoïque $\ce{C6H5COOH_{(aq)}}$ et l'eau.
    \item Calculer la valeur du $\mathrm{pK_{A}}$ du couple
          $\ce{C6H5COOH_{(aq)}\ttslash{}C6H5COO^{-}_{(aq)}}$.
    \item Déterminer, en justifiant votre réponse, l'espèce du couple
          $\ce{C6H5COOH_{(aq)}\ttslash{}C6H5COO^{-}_{(aq)}}$ qui prédomine dans
          la solution $(S_{0})$.
    \item Pour connaître la valeur de la masse $m$ d'acide pur présent dans les
          cristaux préparés, le chimiste a dosé le volume
          $V_{A}=\qty{10,0}{\mL}$ de la solution ($S_{0}$) par une solution
          aqueuse d'hydroxyde de sodium $\ce{Na^{+}_{(aq)} + HO^{-}_{(aq)}}$ de
          concentration molaire $C_{B}=\qty{1,0e-2}{\M}$. Le volume ajouté à
          l'équivalence est $V_{B,E}=\qty{18,0}{\mL}$.
          \begin{enumerate}
              \item Écrire l'équation de la réaction qui se produit entre
                    l'acide benzoïque $\ce{C6H5COOH_{(aq)}}$ et les ions
                    hydroxyde $\ce{HO^{-}_{(aq)}}$ considérée comme totale.
              \item Calculer la valeur de la concentration molaire $C_{A}$ de la
                    solution ($S_{0}$) préparée.
              \item En déduire la valeur de la masse $m$ d'acide benzoïque pur
                    présent dans de la solution ($S_{0}$) de volume $V_{0}$.
              \item Déterminer la valeur du pourcentage $p$ d'acide benzoïque
                    pur contenu dans les cristaux préparés par le chimiste.
          \end{enumerate}
\end{enumerate}

